%
% ferramentas.tex
% Capítulo 3 — Ferramentas, Tecnologias e Utilização de IA
%
\chapter{Ferramentas, Tecnologias e Utilização de Inteligência Artificial}
\label{cha:ferramentas}

\section{Tecnologias e ferramentas utilizadas}
\label{sec:tecnologias-utilizadas}

Descrever as principais tecnologias, linguagens, frameworks, bibliotecas e plataformas utilizadas no desenvolvimento do Play4Change (p.ex. linguagens de programação, motores de jogo/frameworks de front-end e back-end, bases de dados, serviços cloud, ferramentas de design, etc.), justificando as escolhas efetuadas.

\section{Metodologias e processos de trabalho}
\label{sec:metodologias-processos}

Descrever a metodologia de desenvolvimento adotada (p.ex. Scrum/Kanban, controlo de versões, integração contínua, gestão de tarefas) e as ferramentas de apoio utilizadas (p.ex. Git/GitHub, Trello/Jira, Figma).

\section{Utilização de Inteligência Artificial}
\label{sec:utilizacao-ia}

Descrever, de forma transparente, em que fases e de que forma foram utilizadas ferramentas de Inteligência Artificial (p.ex. assistentes de programação, geração/revisão de texto, geração de imagens) ao longo do desenvolvimento deste trabalho, incluindo:

\begin{itemize}
	\item Quais as ferramentas/modelos de IA utilizados;
	\item Em que tarefas foram aplicadas (p.ex. apoio à escrita, geração de código, depuração, criação de conteúdos gráficos, pesquisa);
	\item De que forma o trabalho produzido com apoio de IA foi revisto, validado e integrado pelo autor;
	\item Considerações éticas e limitações identificadas no recurso a estas ferramentas.
\end{itemize}
